\clearpage
\hypertarget{native-479-lower}{}
\begin{center}
\rule{0.55\linewidth}{0.4pt}\\[2pt]
Astronomische Nachrichten. Vol. XXVI. No. 613. pp. 193..208. 14 October 1847.\\[2pt]
\rule{0.55\linewidth}{0.4pt}
\end{center}
\begin{center}[Meridian-circle observations of \emph{Iris} by \textsc{Gauss}.]\end{center}
\begin{flushright}Göttingen, 5 Sept. 1847.\end{flushright}

Mr \textsc{Gould} has just sent me the elements derived from my observations of Iris on 21 and 28 August, together with the Cambridge observation of the 14th, and I hasten to pass them on to you. --- --- --- ---

Last night I made another observation of Iris with the meridian circle.

\begin{center}
\setlength{\tabcolsep}{4pt}
\begin{tabular}{l|r|r|r}
 & Göttingen mean time & R.A. & Declination \\\hline
1847 Sept. 4 & $8^{\rm u}\ 51^{\rm m}\ 12^{\rm s}0$ & $296^\circ\ 7'\ 14''1$ & $-14^\circ\ 6'\ 3''8$ \\
\hfill 6 & $8\ 42\ 54.5$ & $296\ 0\ 46.5$ & $-14\ 8\ 56.1$ \\
\hfill 10 & $8\ 26\ 43.6$ & $295\ 53\ 59.8$ & $-14\ 13\ 59.5$ \\
\end{tabular}
\end{center}

The planet appeared fainter in this observation than in any earlier one. The wires remained quite indistinct if the planet was to be visible. The state of the air may have had much to do with this, for although the sky was completely clear of clouds, all the stars appeared (to the naked eye) in a dim, hazy light.

\begin{center}
\begin{tabular}{l|r|r|r}
Sept. 12 & $8^{\rm u}\ 18^{\rm m}\ 50^{\rm s}5$ & $295^\circ\ 53'\ 37''8$ & $-14^\circ\ 16'\ 16''3$ \\
\end{tabular}
\end{center}

The air was more favourable that evening than on the 10th. On the 13th, however, Iris was just visible only in the completely darkened field, and observation was impossible. The preceding star of magnitude 7. 8., about $1^{\rm m}\ 40^{\rm s}$ ahead and on nearly the same parallel, could still be observed with only scant illumination.

\begin{center}
\begin{tabular}{l|r|r|r}
Sept. 19 & $7^{\rm u}\ 52^{\rm m}\ 19^{\rm s}9$ & $296^\circ\ 8'\ 51''1$ & $-14^\circ\ 22'\ 30''0$ \\
\end{tabular}
\end{center}
